% ============================================================
% RELATED WORK — Full version for NeurIPS 2026
% ============================================================

\section{Related Work}
\label{sec:related}

\paragraph{Classical Granger causality.}
Granger~\citep{granger1969investigating} introduced predictive causality
for bivariate time series, later extended to the multivariate setting
via vector autoregressive (VAR) models~\citep{lutkepohl2005new}. The
framework's key limitation is the assumption of linearity and stationarity.
Nonlinear extensions include kernel Granger
causality~\citep{marinazzo2008kernel} and transfer entropy, but these
remain computationally expensive for multivariate systems and do not
naturally produce causal graphs.

\paragraph{Constraint-based and hybrid methods.}
PCMCI~\citep{runge2019pcmci} addresses the challenge of autocorrelation
in time series causal discovery by combining the PC
algorithm~\citep{spirtes2000causation} with momentary conditional
independence (MCI) tests. It handles contemporaneous confounding better
than bivariate Granger tests but assumes stationarity and scales poorly
with the number of variables and lags. PCMCI has become the standard
baseline in climate science applications, recovering known teleconnections
from reanalysis data. DYNOTEARS~\citep{pamfil2020dynotears} extends the
continuous DAG learning framework NOTEARS to time series, optimizing a
continuous acyclicity constraint. These methods require explicit
statistical testing or optimization over the space of DAGs.

\paragraph{Neural Granger causality.}
Tank et al.~\citep{tank2021neural} proposed component-wise MLPs (cMLP)
and LSTMs (cLSTM) for neural Granger causality, where each output variable
has a dedicated network and group-lasso regularization on input weights
enforces sparsity. When the $\ell_{2,1}$-norm of input weights for variable
$j$ in the network predicting variable $i$ is driven to zero, $X_j$ is
deemed non-causal for $X_i$. This approach handles nonlinear relationships
but requires a causal-specific training objective (prediction loss +
group-lasso penalty), careful tuning of the regularization strength
$\lambda$, and separate training of $K$ networks.

Attention-based methods~\citep{nauta2019causal} use learned attention
weights as proxies for causal strength, but attention weights are known
to be unreliable indicators of feature importance. More recently,
economy SRUs~\citep{khanna2019economy} and self-explaining neural
networks~\citep{marcinkevics2021interpretable} have been proposed for
interpretable Granger causality, while amortized causal
discovery~\citep{lowe2022amortized} learns to infer causal graphs
across datasets. Our approach differs fundamentally: we use
\emph{projection weights} that directly mediate information flow, not
attention-like scores that merely correlate with it.

\paragraph{State space models and interpretability.}
The Mamba architecture~\citep{gu2023mamba} introduced selective state
spaces with data-dependent $\bB$, $\bC$ matrices, achieving competitive
performance with transformers at linear computational cost. The structured
state space duality~\citep{dao2024transformers} showed that SSMs and
attention are mathematically dual, with the $\bB \cdot \bC^\top$ product
acting as an implicit attention matrix. Ali et al.~\citep{ali2024hidden}
explicitly visualized this ``hidden attention'' of Mamba, showing that
SSM internals carry interpretable structure. Le Mercier et
al.~\citep{lemercier2026hispa} demonstrated that hidden state structure
is meaningful for security applications. Our work takes a complementary
direction: rather than interpreting what SSMs \emph{attend to}, we show
that what they \emph{learn to route} (via weight matrices) encodes causal
structure.

\paragraph{Causal representation learning.}
Sch\"olkopf et al.~\citep{scholkopf2021toward} advocate learning
representations that respect causal structure, arguing that causal
variables should be the building blocks of machine learning. Our work
provides a concrete instantiation of this vision: a standard prediction
model (Mamba), trained without any causal objective, naturally learns
representations whose structure reflects the underlying causal graph.
This is a \emph{descriptive} rather than \emph{prescriptive} result---we
do not design the model to be causal, but discover that it is.

\paragraph{Positioning of our work.}
\methodname{} is most closely related to neural Granger causality
(cLSTM~\citep{tank2021neural}) but differs in three fundamental ways:
(1)~we require no causal-specific training objective;
(2)~we use a single shared model rather than $K$ separate networks; and
(3)~our method works post-hoc on any pre-trained SSM, enabling causal
analysis of existing models without retraining. Relative to SSM
interpretability work, we are the first to connect SSM parameters to
\emph{causal} rather than merely \emph{attentional} structure.
